\section[Finite Metric Separation for 1-NN]{Finite Metric Separation for $1$-NN}
\label{sec:finite-separation}

We now come to the main theorem-level witness of the draft. The statement is
small, explicit, and completely deterministic.

\begin{proposition}[Pointwise LOO does not control pointwise replace-one for deterministic $1$-NN]
\label{prop:two-point-separation}
There exists a finite graph metric space, an ordered labeled sample $S$, and a
query-label pair $(x,y)$ such that
\[
\Delta_{\loo}(S,i)=0 \quad\text{for every sample index } i,
\]
while for some index $j$ and some admissible replacement point $z$,
\[
\Delta_{\rep}(S,j,z,x,y)=1.
\]
\end{proposition}

\begin{proof}
Let the metric space contain two vertices $a,b$ joined by one edge, so
$d(a,b)=1$. Consider the ordered sample
\[
S=((a,0),(b,0)).
\]
Use deterministic $1$-NN with the fixed tie-breaking conventions from
\Cref{sec:definitions}.

For LOO, delete the first occurrence. The remaining sample is $((b,0))$, whose
$1$-NN prediction at the deleted occurrence's location $a$ is still $0$ because
the sole remaining neighbor has label $0$. Hence the loss at $(a,0)$ is
unchanged. The same argument applies when deleting the second occurrence and
evaluating at $(b,0)$. Therefore $\Delta_{\loo}(S,1)=\Delta_{\loo}(S,2)=0$.

Now evaluate at the fixed query-label pair $(x,y)=(a,0)$. Replace the first
occurrence $(a,0)$ by $(a,1)$, obtaining
\[
S^{1\leftarrow (a,1)}=((a,1),(b,0)).
\]
At query $a$, the original sample predicts $0$ because $(a,0)$ is the unique
nearest occurrence. In the perturbed sample, the unique nearest occurrence is
now $(a,1)$, so the prediction becomes $1$. Hence
$\Delta_{\rep}(S,1,(a,1),a,0)=1$.
\end{proof}

\Cref{fig:minimal-witness} gives the same witness visually. What matters is not
that the example is exotic, but that it is minimal enough to inspect
line-by-line. The phenomenon already appears before any asymptotics, averages,
or large combinatorial constructions enter the story.

\begin{table}[t]
\centering
\small
\begin{tabular}{lllllll}
\toprule
Vertices & Edges & Sample & k & max LOO & max replace-one & Status \\
\midrule
2 & 1 & order=[0, 1]; labels=[0, 0] & 1 & 0 & 1 & explicit witness; certificate-level minimality only \\
2 & 1 & order=[0, 1]; labels=[1, 1] & 1 & 0 & 1 & explicit witness; certificate-level minimality only \\
2 & 1 & order=[1, 0]; labels=[0, 0] & 1 & 0 & 1 & explicit witness; certificate-level minimality only \\
2 & 1 & order=[1, 0]; labels=[1, 1] & 1 & 0 & 1 & explicit witness; certificate-level minimality only \\
\bottomrule
\end{tabular}
\caption*{\footnotesize The witness row is hand-checkable, while minimality over the search space remains computational evidence.}
\end{table}


The repository also contains a finite-search pipeline that looks for smaller or
structurally simpler witnesses under various restrictions. The current output is
best interpreted as follows. The witness in \Cref{prop:two-point-separation} is
theorem-level because it is explicit and hand-checkable. Claims that no witness
exists below a certain search frontier are \emph{not} theorem-level; they are
computational certificates tied to the enumerated search space.

\begin{table}[t]
\centering
\small
\begin{tabular}{lllll}
\toprule
Artifact & Observed min vertices & No-solution range & Hash & Status \\
\midrule
TASK-007 & 2 & [1] & 5b7a8925688a... & computational evidence only \\
TASK-008 & 2 & [1] & a52aa8a77a9a... & computational evidence only \\
\bottomrule
\end{tabular}
\caption*{\footnotesize Certificate tables summarize finite search outputs and do not upgrade them to proof.}
\end{table}


This distinction matters. A paper can responsibly say ``here is a minimal
explicit witness we know'' and separately say ``here is a reproducible finite
certificate supporting minimality within the searched regime.'' It should not
blur those two statements together. That boundary is part of the contribution,
not an embarrassment to hide.
